% chapters/03_verwandte_arbeiten.tex
%
% Überarbeitete Fassung vom 16.10.2026.
% Alle Angaben zu zitierten Arbeiten wurden gegen die Volltexte geprüft
% (arXiv-HTML/PDF, ACL Anthology, ICLR-Proceedings, PMC). Angaben zur
% eigenen Arbeit folgen Kapitel 1 und 4.
% Benötigt in der Präambel: booktabs, array (über tabularx geladen).

\chapter{Verwandte Arbeiten}
\label{chap:verwandte_arbeiten}

Dieses Kapitel ordnet die vorliegende Arbeit in den Forschungsstand zur
Named Entity Recognition (NER) mit Encoder-Modellen und generativen Large
Language Models (LLMs) ein.
Berücksichtigt werden Arbeiten, die mindestens eines von vier Kriterien
erfüllen: Sie vergleichen Encoder und generative LLMs auf derselben
NER-Aufgabe, sie untersuchen Ausgabeformate oder Anpassungsverfahren für
LLM-basierte NER, sie verwenden MultiNERD oder CrossNER, oder sie
variieren die Modellgröße.
Begutachtete Fassungen werden bevorzugt; Preprints werden einbezogen,
wenn keine begutachtete Fassung vorliegt und die Arbeit unmittelbar
einschlägig ist, darunter zwei Arbeiten aus dem Jahr
2026~\cite{zhan2026generativeNER,semin2026spanLabeling}.
Die Auswahl beruht auf einer narrativen Recherche in Google Scholar,
arXiv und der ACL Anthology (2019 bis 2025) mit Stichwörtern wie
\enquote{LLM named entity recognition}, \enquote{generative NER},
\enquote{instruction tuning NER}, \enquote{LoRA/QLoRA NER},
\enquote{DeBERTa NER}, \enquote{MultiNERD} und \enquote{CrossNER}.
Ergänzend wurde entlang der Zitationen der Schlüsselarbeiten per
Schneeballsuche gesucht.
Trefferzahlen wurden nicht protokolliert; die Auswahl beansprucht keine
systematische Vollständigkeit.
Grundlagen zu Architekturen, LoRA und QLoRA, Skalierungsgesetzen,
Datensätzen und Metriken wurden in Kapitel~\ref{chap:grundlagen}
eingeführt und werden hier nicht wiederholt.

Abschnitt~\ref{sec:encoder_ner_related} behandelt Encoder-basierte
Ansätze, einschließlich der MultiNERD-Originalarbeit und solcher Encoder,
die mit LLM-erzeugten Trainingsdaten als Generalisten eingesetzt werden.
Abschnitt~\ref{sec:llm_ner_related} gliedert LLM-basierte Ansätze nach
Ausgabestrategie und Anpassungsart und begründet daraus die in dieser
Arbeit gewählte Ausgabestrategie.
Abschnitt~\ref{sec:peft_ner_related} stellt Arbeiten vor, die mit
Parameter-Efficient Fine-Tuning (PEFT) angepasste LLMs direkt mit
Encodern vergleichen.
Abschnitt~\ref{sec:model_size_related} fasst Befunde zum Einfluss der
Modellgröße zusammen.
Abschnitt~\ref{sec:research_gap} leitet daraus die Forschungslücke ab und
ordnet sie den Forschungsfragen aus Abschnitt~\ref{sec:problem_ziel} zu.
Eigene Ergebnisse werden nicht vorweggenommen.

\section{Encoder-basierte Ansätze für NER}
\label{sec:encoder_ner_related}

Encoder mit Klassifikationskopf bilden seit BERT den Standard für
überwachte NER (Abschnitt~\ref{sec:encoder_ner}).
Für diese Arbeit sind drei Fragen leitend: Wie ist die Datenbasis
MultiNERD entstanden und welche Referenzwerte liegen für sie vor?
Wie leistungsfähig sind feinjustierte Encoder bei feingranularer NER im
Vergleich zu per Prompting genutzten LLMs?
Und wie weit tragen Encoder, wenn sie ohne aufgabenspezifische
Annotationen eingesetzt werden?

\textcite{tedeschi2022multinerd} stellen mit MultiNERD einen Datensatz
vor, der zehn Sprachen, die Genres Wikipedia und Wikinews sowie
15 Entitätstypen abdeckt; der englische Teil umfasst rund 164\,000 Sätze.
Die Annotation erfolgt automatisch: Verlinkte Textstellen in Wikipedia-
und Wikinews-Artikeln werden über BabelNet-Synsets einem Entitätstyp
zugeordnet, auf nicht verlinkte Vorkommen übertragen und anschließend
durch ein iteratives Verfahren bereinigt, das Sätze verwirft, bei denen
ein auf den Silberdaten trainiertes Modell (mBERT mit BiLSTM und CRF) der
wissensbasierten Annotation widerspricht.
Dieses mBERT-Modell ist zugleich die einzige in der Arbeit trainierte
Architektur.
Evaluiert wird nicht auf den eigenen Splits, sondern einerseits nach
Abbildung auf die vier groben CoNLL-Typen auf Gold-Benchmarks wie
CoNLL-2003, andererseits auf 1\,000 englischen Sätzen, deren automatische
Annotation manuell geprüft wurde.
Auf dieser Stichprobe stimmt die Silberreferenz auf Tokenebene mit
einem F$_1$ von 0{,}85 mit der manuellen Korrektur überein; PER, ORG,
LOC, CEL, DIS, EVE und MEDIA zählen zu den zuverlässig annotierten
Klassen, während ANIM, BIO, FOOD und PLANT häufig verwechselt oder nicht
verlinkt werden.
Für die offiziellen Splits, die diese Arbeit verwendet
(Abschnitt~\ref{sec:datenauswahl}), berichtet die Originalarbeit somit
keine Referenzwerte.
Veröffentlichte Werte auf MultiNERD finden sich in der hier gesichteten
Literatur nur in Benchmarks instruktionsgetunter LLMs
(Abschnitt~\ref{sec:llm_ner_related}) und in der unten behandelten
GLiNER-Arbeit; sie verwenden eigene Datenaufbereitungen, oder ihr
Testumfang lässt sich nicht auf die offiziellen Splits abbilden.

Wie leistungsfähig Encoder bei feingranularer NER sind, zeigt der
MultiCoNER-II-Wettbewerb mit 33 feingranularen Typen in zwölf Sprachen.
\textcite{tan2023damo} gewinnen dort mit U-RaNER neun der 13 Tracks.
Das System kombiniert XLM-RoBERTa-large mit einer CRF-Schicht und einem
Retrieval-Modul, das zu jedem Eingabesatz Textpassagen aus Wikipedia
sowie Entitäten und Typinformationen aus Wikidata abruft und dem Encoder
als zusätzlichen Kontext mitgibt; auf dem englischen Track erreicht es
einen Makro-F$_1$ von 85{,}53.
Ergänzend prompten die Autoren ChatGPT (\texttt{gpt-3.5-turbo}) in drei
Varianten auf dem Validierungssplit.
Selbst die beste Variante mit Demonstrationsbeispielen erreicht im Mittel
über die Sprachen nur einen Mikro-F$_1$ von 14{,}78, gegenüber 92{,}77
für das eigene System auf demselben Split~\cite[Tab.~4]{tan2023damo}.
Der Befund bezieht sich auf ein Modell ohne jede aufgabenspezifische
Anpassung und auf ein sehr feines Typinventar; er zeigt aber, wie groß
der Abstand zwischen reinem Prompting und überwachten Encodern bei
feingranularer NER ausfallen kann.

Eine breitere Einordnung liefert die Vergleichsstudie von
\textcite{keraghel2024recent}, die als Teil eines Surveys zehn
NER-Datensätze aus allgemeinen und spezialisierten Domänen mit mehreren
Frameworks und Modellen auswertet, darunter BERT, XLM-RoBERTa-large,
DeBERTa-v3-base und DeBERTa-v3-large, ein feinjustiertes GLiNER-Modell
sowie GPT-4o mit Beispielentitäten im Prompt.
MultiNERD ist nicht enthalten.
Unter den Transformer-Modellen erweisen sich die DeBERTa-Varianten als am
effektivsten; nach einem Nemenyi-Test teilen sich DeBERTa-v3-large und
das feinjustierte GLiNER-Modell die besten Ränge, ohne sich signifikant
zu unterscheiden.
GPT-4o bleibt auf neun der zehn Datensätze deutlich hinter den
feinjustierten Encodern zurück, auf CoNLL-2003 etwa mit einem
Makro-F$_1$ von 65{,}17 gegenüber 90{,}98 für
DeBERTa-v3-large~\cite[Tab.~5]{keraghel2024recent}.
Auf WNUT2017 liegt GPT-4o mit 44{,}73 nahezu gleichauf mit
DeBERTa-v3-large (44{,}81), bleibt aber auch dort hinter dem
feinjustierten GLiNER-Modell zurück.
Die Autoren folgern, dass LLMs allein für NER nur eingeschränkt geeignet
sind und eine Kombination mit feinjustierten Modellen vielversprechender
erscheint.
Die Studie stützt die Wahl von DeBERTa-v3 als Encoder-Referenz in dieser
Arbeit (Abschnitt~\ref{sec:encoder-auswahl}); sie liegt bislang nur als
Preprint vor.

Eine zweite Entwicklungslinie nutzt Encoder als Generalisten, die
Entitätstypen zur Laufzeit aus natürlichsprachlichen Bezeichnungen
ableiten.
\textcite{zaratiana2024gliner} stellen mit GLiNER ein Modell auf Basis
von DeBERTa-v3 vor, das Typbezeichnungen und Eingabetext in einer
gemeinsamen Sequenz verarbeitet, Kandidatenspans bis zu einer
Maximallänge repräsentiert und jeden Span über ein Skalarprodukt mit den
Typrepräsentationen bewertet.
Trainiert wird auf Pile-NER, den von ChatGPT erzeugten Annotationen, die
auch UniversalNER verwendet (Abschnitt~\ref{sec:llm_ner_related}).
Im Zero-Shot-Vergleich auf sieben Datensätzen außerhalb der
Trainingsverteilung, darunter die fünf CrossNER-Domänen, erreicht die
mittlere Variante mit rund 90~Mio. Parametern im Mittel einen F$_1$ von
55{,}4 und damit das Niveau des etwa 140-mal größeren UniNER-13B; die
größte Variante liegt im Mittel rund fünf Punkte darüber, und selbst die
kleinste Variante mit 50~Mio. Parametern übertrifft ChatGPT.
Das ebenfalls verglichene GoLLIE-7B ordnet sich zwischen der mittleren
und der größten Variante ein.
Inferenzzeiten messen die Autoren nicht; sie verweisen auf die parallele
Span-Klassifikation gegenüber der sequenziellen Generierung von LLMs und
auf die Lauffähigkeit auf CPUs.

\textcite{bogdanov2024nuner} verfolgen mit NuNER einen verwandten Weg,
setzen aber auf ein Pre-Training für die anschließende Feinjustierung
statt auf das Zero-Shot-Regime.
Ein RoBERTa-base-Modell wird kontrastiv auf 4{,}38~Mio.
Entitätsannotationen vortrainiert, die GPT-3.5 für Sätze aus dem
C4-Korpus über rund 200\,000 verschiedene Konzepte erzeugt hat.
Zwei Vergleiche mit LLMs sind dabei zu unterscheiden.
Gegenüber UniNER-7B mit 56-mal mehr Parametern zeigt NuNER im
Few-Shot-Fine-Tuning auf den vollständigen Testsets von vier Datensätzen
(Mikro-F$_1$ auf Entitätsebene) ähnliche
Transferfähigkeiten~\cite[Tab.~2]{bogdanov2024nuner}.
GPT-3.5 und GPT-4 per In-Context-Learning werden nur auf zwei
Datensätzen mit auf 1\,000 Beispiele verkleinerten Testsets und
Makro-F$_1$ auf Tokenebene verglichen; NuNER übertrifft dort GPT-3.5 bei
größeren Trainingsmengen, während die Autoren das Verhältnis zu GPT-4
selbst als unklar bezeichnen~\cite[Abschn.~6]{bogdanov2024nuner}.

Zusammengenommen zeigen diese Arbeiten zweierlei.
Erstens bleiben vollständig feinjustierte Encoder im überwachten Setting
die Referenz, auch bei feingranularen Typinventaren, und per Prompting
genutzte LLMs ohne Anpassung liegen in den hier berichteten Vergleichen deutlich
zurück~\cite{tan2023damo,keraghel2024recent}.
Zweitens lässt sich Wissen aus LLMs über erzeugte Trainingsdaten in
Encoder übertragen, sodass Modelle mit einem Bruchteil der Parameter mit
LLMs im Bereich von 7 bis 13~Mrd. Parametern
gleichziehen~\cite{zaratiana2024gliner,bogdanov2024nuner}.
Diese Gleichstände beruhen jedoch auf Zero- und Few-Shot-Benchmarks, bei
GLiNER unter anderem auf den domänenspezifischen Typinventaren von
CrossNER; keine der Arbeiten stellt vollständig feinjustierte Encoder
und feinjustierte generative LLMs in derselben Pipeline auf den
offiziellen MultiNERD-Splits mit gemeinsamer Ressourcenmessung
gegenüber.
Für MultiNERD kommt hinzu, dass die Originalarbeit keine Referenzwerte
für die offiziellen Splits liefert und die Testreferenz selbst ein Silver
Standard ist.
Vergleiche auf diesem Datensatz müssen daher innerhalb einer gemeinsamen
Pipeline erfolgen und absolute Werte als Übereinstimmung mit einer
automatisch erzeugten Referenz lesen (Abschnitt~\ref{sec:datenauswahl}).

\section{LLM-basierte Ansätze für NER}
\label{sec:llm_ner_related}

Wird NER als Textgenerierung formuliert, sind zwei Entwurfsentscheidungen
zu treffen, die in der Literatur getrennt untersucht werden: das
Ausgabeformat, in dem das Modell Entitäten benennt, und die Art der
Anpassung, von reinem Prompting über Destillation bis zum Fine-Tuning.
Dieser Abschnitt folgt dieser Zweiteilung.
Zunächst werden Ausgabestrategien und ihre Wirkung auf Erkennungsleistung
und Ausgabelänge dargestellt, woraus die Begründung der in dieser Arbeit
gewählten Strategie folgt.
Anschließend werden Ansätze ohne aufgabenspezifisches Training und
Ansätze mit Destillation oder Fine-Tuning vorgestellt.

\textcite{semin2026spanLabeling} systematisieren die Ausgabestrategien
für Span-Labeling mit LLMs in drei Familien.
Beim \emph{Tagging} wird der Eingabetext mit Markierungen reproduziert,
sodass die Ausgabelänge linear mit der Eingabelänge wächst.
Beim \emph{Indexing} gibt das Modell für jeden Span numerische Start- und
Endindizes aus, sodass die Ausgabelänge nur von der Zahl der Spans
abhängt.
Beim \emph{Matching} erzeugt das Modell den Text der Spans, der
anschließend per Zeichenkettenabgleich in der Eingabe lokalisiert wird.
Die Autoren vergleichen die Strategien mit vier per Prompting genutzten
offenen Modellen zwischen 8 und 70~Mrd. Parametern unter anderem auf dem
mehrsprachigen UNER-Benchmark mit den Typen Person, Organisation und
Ort.
Dieser Benchmark ist der Datensatz Universal NER von
\textcite{mayhew2024uner}, nicht das UniversalNER-Modell von
\textcite{zhou2024universalner}.
Tagging erweist sich als robusteste, aber tokenintensivste Strategie;
Matching ist konkurrenzfähig und tokeneffizient; Indexing mit
Zeichenindizes ist nur dann konkurrenzfähig, wenn die Eingabe selbst mit
Zeichenindizes angereichert wird, weil die Modelle Zeichenpositionen
sonst nur indirekt schätzen können.
Mit LogitMatch schlagen die Autoren zudem ein Constrained-Decoding-Verfahren
vor, das Matching-Ausgaben auf gültige Eingabespans beschränkt und so
Fehler beim Zeichenkettenabgleich beseitigt.
Eine auf \textcite[Tab.~9]{ramponi2025fallacy} zurückgehende
Indexing-Variante mit nummerierten Wörtern statt Zeichenpositionen wird
beschrieben, aber nicht evaluiert.
Die Arbeit liegt bislang nur als Preprint vor.

Dass das Format auch bei feinjustierten Modellen ein Faktor erster
Ordnung ist, zeigen \textcite{zhan2026generativeNER}.
Sie trainieren acht offene LLMs mit LoRA auf fünf Ausgabeformaten: zwei
Inline-Formate, die Entitäten mit Klammern oder XML-Tags im
reproduzierten Satz markieren, und drei JSON-Formate, die Entitäten nach
Typ gruppieren, mit einem Vorkommensindex versehen oder mit Zeichenoffsets
angeben.
Die Inline-Formate erreichen über alle Modelle die besten Mittelwerte,
auf den flachen Datensätzen rund 92 F$_1$, während JSON mit Zeichenoffsets
dort im Mittel auf rund 54 fällt; die Autoren führen dies darauf zurück,
dass das Berechnen von Zeichenpositionen der generativen Arbeitsweise der
Modelle schlecht entspricht.
\textcite{dagdelen2024structuredIE} berichten für die Extraktion aus
materialwissenschaftlichen Texten mit feinjustiertem GPT-3 und einem
LoRA-angepassten Llama-2-Modell, dass JSON-Ausgaben in rund 99\,\% der
Fälle parsebar sind, die Formatierung aber nicht erzwungen wird und JSON
mehr Tokens benötigt als satzförmige Schemata.
Auf die Kosten der Ausgabelänge weisen \textcite{lu2024padellm} hin: Die
Latenz generativer NER entsteht vor allem durch die sequenzielle
Dekodierung langer Ausgaben, und Inline-Formate verlängern die Ausgabe,
weil sie den gesamten Eingabetext wiederholen.
Ihr paralleles Decoding, das Erwähnungen je Typ gleichzeitig statt
nacheinander erzeugt, verkürzt die Inferenz um das 1{,}76- bis
10{,}22-Fache bei unveränderter Erkennungsleistung.

Vor diesem Hintergrund lässt sich die in Abschnitt~\ref{sec:prompt-design}
beschriebene Ausgabestrategie dieser Arbeit einordnen.
Das Modell erhält eine nummerierte Tokenliste und gibt ein JSON-Array mit
Start- und Endindex, Oberflächentext und Typ je Entität aus.
Das Format gehört zur Indexing-Familie, allerdings in der angereicherten
Variante auf Wortebene: Die nummerierten Eingabetokens und die
Wortindizes in der Ausgabe folgen \textcite{ramponi2025fallacy}, die
Einordnung in die drei Familien \textcite{semin2026spanLabeling};
JSON-Schema und Nachverarbeitung sind eigene Umsetzung dieser Arbeit.
Daraus ergeben sich drei erwartete Vorteile.
Erstens stehen die Indizes explizit in der Eingabe, sodass das Modell
keine Zeichen zählen muss, was den von \textcite{zhan2026generativeNER}
beobachteten Einbruch bei Zeichenoffsets vermeiden soll.
Zweitens hängt die Ausgabelänge von der Zahl der Entitäten, der Länge
ihrer Spantexte und der Serialisierung ab; anders als bei Tagging- und
Inline-Formaten entfällt lediglich die Reproduktion der
Nicht-Entitätstokens.
Drittens erlauben die Indizes eine eindeutige Rückabbildung auch bei
mehrfach vorkommenden Oberflächenformen, die beim Matching eine
Fehlerquelle sind.
Das redundante Textfeld ermöglicht die Konsistenzprüfung in der
Nachverarbeitung (Abschnitt~\ref{sec:postprocessing}).
Einschränkend ist festzuhalten, dass \textcite{semin2026spanLabeling} die
Wortindex-Variante nicht evaluiert haben und ihre Befunde auf
per Prompting genutzten, nicht feinjustierten Modellen beruhen; ob die Vorteile der
angereicherten Indizierung auch im QLoRA-Regime und bei 15 Typen tragen,
ist eine empirische Frage dieser Arbeit.
Auf Constrained Decoding wird verzichtet, damit die Formtreue der
Ausgaben selbst messbar bleibt (Abschnitt~\ref{sec:robustheit-ausgaben}).

Die frühen LLM-Ansätze für NER lassen die Gewichte des Modells
unverändert und passen es nicht an die Aufgabe an.
\textcite{wang2025gptner} formulieren mit GPT-NER die Aufgabe als
Generierung des Eingabesatzes, in dem Entitäten eines Typs durch die
Marker \texttt{@@} und \texttt{\#\#} eingefasst sind, und wählen
Demonstrationsbeispiele über eine Nachbarschaftssuche auf
Entitätsrepräsentationen aus.
Diese Repräsentationen liefert in der stärksten Variante ein
feinjustierter NER-Tagger~\cite[Abschn.~3.2.2]{wang2025gptner}.
Das verwendete GPT-3-Modell neigt dazu, auch solche Eingaben
überkonfident mit Entitäten zu versehen, die gar keine Entität des
abgefragten Typs enthalten; die Autoren bezeichnen diese Fälle als
\emph{NULL inputs} und ergänzen eine Selbstverifikation, in der das
Modell jede extrahierte Entität mit Ja oder Nein bestätigt.
Auf CoNLL-2003 erreicht GPT-NER so einen F$_1$ von 90{,}91 gegenüber
92{,}8 für einen feinjustierten BERT-Tagger~\cite[Tab.~1]{wang2025gptner};
bei verschachtelten Datensätzen ist der Abstand größer.
\textcite{xie2023zeroshot} untersuchen Zero-Shot-NER mit ChatGPT und
zerlegen die Aufgabe in eine Abfrage je Entitätstyp, ergänzt um
syntaktische Hinweise im Prompt und ein zweistufiges Mehrheitsvotum über
mehrere Antworten.
Die Zerlegung verbessert die Ergebnisse auf allen sieben untersuchten
Benchmarks, vervielfacht aber die Zahl der Modellaufrufe je Satz; die in
den Tabellen der Arbeit mitgeführten überwachten Bestwerte liegen
weiterhin weit über den Zero-Shot-Ergebnissen.

\textcite{ashok2023promptner} erweitern das Prompting mit PromptNER um
textuelle Typdefinitionen und verlangen vom Modell eine Kandidatenliste
mit Begründungen, warum ein Kandidat zur Definition passt oder nicht.
Mit GPT-4 und wenigen Demonstrationsbeispielen verbessert PromptNER in
der zweiten Fassung der Arbeit die besten bisherigen
Few-Shot-Ergebnisse auf CoNLL-2003 um rund vier, auf GENIA um rund neun
und auf FewNERD um rund vier F$_1$-Punkte.
Im Cross-Domain-Setting mit nur zwei Beispielen aus der Zieldomäne
erreicht es auf drei der fünf CrossNER-Domänen neue Bestwerte, etwa
84{,}26 auf \emph{Music}, obwohl weniger als 2\,\% der verfügbaren
Trainingsdaten verwendet werden; in den Domänen \emph{AI} und
\emph{Science} bleiben die bisherigen Verfahren vorn.

Einen anderen Weg geht CodeIE von \textcite{li2023codeie}: NER und
Relationsextraktion werden als Vervollständigung von Python-Code
formuliert und mit dem Code-Modell Codex gelöst.
Im Few-Shot-Setting übertrifft CodeIE sowohl natürlichsprachliche Prompts
als auch feinjustierte UIE-Modelle, die nur dieselben wenigen Beispiele
sehen; auf CoNLL-2003 steht einem Wert von 82{,}32 mit fünf Beispielen je
Typ jedoch ein vollständig überwachter UIE-Wert von 92{,}99 gegenüber.

Eine zweite Gruppe passt offene LLMs gezielt an NER an.
\textcite{zhou2024universalner} lassen ChatGPT Passagen aus dem
Pile-Korpus mit Entitäten frei gewählter Typen annotieren und trainieren
damit LLaMA-Modelle mit 7 und 13~Mrd. Parametern als UniversalNER.
Auf einem Benchmark aus 43 Datensätzen in neun Domänen übertreffen die
Studentenmodelle im Zero-Shot-Regime ihren Lehrer deutlich, mit
mittleren F$_1$-Werten von 41{,}7 bzw. 43{,}4 gegenüber 34{,}9 für
ChatGPT.
Das Training folgt dem Rezept von Vicuna; Adapterverfahren werden nicht
erwähnt, Ressourcen nicht gemessen.

In einem zusätzlichen überwachten Setting wird UniNER-7B auf den
Trainingsdaten von 20 Datensätzen weiter feinjustiert und erreicht im
Mittel 84{,}78 gegenüber 80{,}09 für BERT-base; die BERT-Werte sind aus
der InstructUIE-Arbeit~\cite{wang2023instructuie} übernommen und nicht neu
gemessen.
Zu diesen 20 Datensätzen gehört auch eine Variante von MultiNERD, auf
der UniNER-7B 93{,}73 und BERT-base 91{,}25
erreichen~\cite[Tab.~2]{zhou2024universalner}.
Da InstructUIE eigene Datenaufbereitungen verwendet, lassen sich diese
Werte nicht auf die offiziellen MultiNERD-Splits übertragen; sie zeigen
aber, dass auf diesem Datensatz sowohl Encoder als auch feinjustierte
LLMs hohe Werte erreichen und der Abstand zwischen beiden klein ist.

\textcite{sainz2024gollie} verfolgen mit GoLLIE das Ziel,
Annotationsrichtlinien explizit nutzbar zu machen.
Code-LLaMA-Modelle mit 7, 13 und 34~Mrd. Parametern werden mit QLoRA auf
Informationsextraktionsaufgaben feinjustiert, deren Schemata samt
Richtlinien als Python-Klassen im Prompt stehen.
Im Zero-Shot-Regime auf nicht im Training enthaltenen Datensätzen hebt
das Richtlinientraining den mittleren F$_1$ gegenüber einer Baseline mit
demselben Backbone ohne Richtlinien (und ohne die zugehörigen
Kandidatenlisten und Regularisierung) um rund 13 Punkte auf 55{,}3 für
das 7B-Modell; die größeren Modelle legen nur wenig zu.
Im überwachten Setting dagegen unterscheiden sich GoLLIE und Baseline
kaum.
Auf den fünf CrossNER-Domänen erreicht GoLLIE-7B Werte zwischen 55{,}5
(\emph{Science}) und 67{,}8 (\emph{Music}) und übertrifft damit das bis
dahin beste Zero-Shot-System InstructUIE in jeder
Domäne~\cite[Tab.~3]{sainz2024gollie}.

Auch MultiNERD ist Teil der Zero-Shot-Evaluation von GoLLIE.
Das 7B-Modell erreicht dort einen F$_1$ von 76{,}0 und liegt damit
deutlich über der Baseline ohne Richtlinien~\cite[Tab.~3]{sainz2024gollie},
wobei acht der 15 Typen in anderer Benennung bereits im Training
vorkamen.
Die Arbeit gibt an, für alle Datensätze die Standardsplits zu verwenden,
und nennt für MultiNERD 32\,908 Testbeispiele, genau das Doppelte der
16\,454 Sätze des englischen Testsplits dieser Arbeit
(Abschnitt~\ref{sec:datenaufteilung}); wie beide Angaben zusammenhängen,
bleibt ungeklärt, sodass der Wert nur eingeschränkt vergleichbar
ist~\cite[Anh.~D.2, Tab.~9]{sainz2024gollie}.

Einen weiteren Weg zur Zero-Shot-Generalisierung geht GNER von
\textcite{ding2024gner}.
Flan-T5-Modelle bis 11~Mrd. Parameter und LLaMA-7B lernen auf Pile-NER,
jedes Eingabewort mit seinem Label zu wiederholen, also auch
Nicht-Entitäten als negative Instanzen mitzuführen, was nach den Autoren
Kontext liefert und Labelgrenzen schärft.
Ein Abgleich über die längste gemeinsame Teilfolge ordnet die Ausgabe
auch bei ausgelassenen oder hinzugefügten Wörtern der Eingabe zu.
Im Zero-Shot-Regime auf den CrossNER- und MIT-Datensätzen berichten die
Autoren einen Gewinn von 9 F$_1$-Punkten gegenüber dem bisherigen Stand
der Technik.
Verglichen wird mit Werten anderer Systeme, darunter GLiNER, die teils
übernommen und teils mit veröffentlichten Checkpoints neu evaluiert
wurden; gemessen werden Inferenzdurchsatz und gesamte GPU-Stunden des
Trainings, aber kein Encoder in derselben Pipeline.
Das Label-je-Wort-Format reproduziert den gesamten Eingabetext und
gehört damit zur Tagging-Familie, anders als das hier verwendete
Indexing-Format.

Über die Einzelarbeiten hinaus ergeben sich drei Beobachtungen.
Erstens ist das Ausgabeformat kein Implementierungsdetail, sondern
bestimmt Erkennungsleistung und Ausgabelänge zugleich; Formate, die den
Eingabetext reproduzieren, sind robust, aber teuer, Formate mit
Zeichenpositionen
fehleranfällig~\cite{semin2026spanLabeling,zhan2026generativeNER,lu2024padellm}.
Zweitens bleiben per Prompting genutzte LLMs ohne Anpassung der Gewichte
auf etablierten Benchmarks hinter überwachten Encodern zurück, selbst
wenn wie bei GPT-NER ein feinjustierter Tagger die Demonstrationen
auswählt; ihre Stärke zeigen sie dort, wo wenige Beispiele vorliegen
oder Domänen
wechseln~\cite{wang2025gptner,xie2023zeroshot,ashok2023promptner,li2023codeie}.
Drittens verringern Destillation und Fine-Tuning diesen Abstand, doch die
Vergleiche mit Encodern stützen sich auf übernommene Werte oder auf
Zero-Shot-Benchmarks~\cite{zhou2024universalner,sainz2024gollie}.
UniversalNER misst keinen Trainings- oder Inferenzaufwand; GoLLIE
dokumentiert Trainingszeit und CO$_2$ je Größe (7B: eine A100,
44{,}5~h; 13B: 79{,}5~h; 34B: zwei A100, 94{,}6~h), aber keine
Inferenzlatenz und keinen
Encodervergleich~\cite[Anh.~D.3, Tab.~10]{sainz2024gollie}.
Für den CrossNER-Ergänzungsversuch dieser Arbeit
(Abschnitt~\ref{sec:crossner}) liefern PromptNER und GoLLIE Orientierung,
aber keine Vergleichswerte: PromptNER sieht zwei Beispiele aus der
Zieldomäne, GoLLIE trainiert auf den Richtlinien vieler Datensätze, und
beide werten die domänenspezifischen Typinventare von CrossNER aus.
Der Versuch dieser Arbeit ist dagegen ein reiner Transfertest ohne
Zielbeispiele, beschränkt auf die Grobtypen PER, ORG und LOC.
Die berichteten Werte zeigen deshalb lediglich, dass der Domänenwechsel
auch für generative Modelle eine erhebliche Hürde bleibt, nicht aber,
welches Niveau auf dem reduzierten Typinventar zu erwarten ist.

\section[PEFT-angepasste LLMs im direkten Vergleich mit Encodern]{PEFT-angepasste LLMs im direkten \mbox{Vergleich} mit Encodern}
\label{sec:peft_ner_related}

LoRA und QLoRA wurden in Abschnitt~\ref{sec:lora} als Verfahren
eingeführt.
Dieser Abschnitt behandelt Arbeiten, die generative LLMs mit solchen
Adaptern an NER anpassen und in derselben Studie mit Encodern
vergleichen.
Soweit die hier gesichtete Literatur reicht, ist diese Kombination noch
selten untersucht; die vier im Folgenden dargestellten Arbeiten decken
allgemeine und domänenspezifische Datensätze ab und unterscheiden sich
darin, welche Ressourcen sie erfassen.

Den breitesten Vergleich auf allgemeinen Benchmarks legen
\textcite{zhan2026generativeNER} vor.
Acht offene LLMs der Familien LLaMA, Qwen, Granite und Gemma mit 3 bis
9~Mrd. Parametern werden mit LoRA auf CoNLL-2003 und OntoNotes 5.0 sowie
auf den verschachtelten Datensätzen ACE2005 und GENIA feinjustiert und
spezialisierten Modellen wie BERT-MRC gegenübergestellt.
Im besten Format erreicht LLaMA3.1-8B auf CoNLL-2003 einen F$_1$ von
93{,}85 gegenüber 93{,}04 für BERT-MRC~\cite[Tab.~3]{zhan2026generativeNER};
auf GENIA liegt das beste LLM 4{,}02 Punkte hinter dem besten
spezialisierten Modell.
Für den Kostenvergleich trainieren die Autoren LLaMA3.1-8B und BERT-MRC
selbst auf einer einzelnen L40-GPU und berichten dort dieselben
F$_1$-Werte wie in der Haupttabelle~\cite[Anh.~G, Tab.~13]{zhan2026generativeNER};
für die übrigen Baselines nennt die Arbeit keine Herkunft der Werte.
Eine Fehleranalyse für LLaMA3.1-8B und BERT-MRC zeigt unterschiedliche
Profile: Der Encoder übersieht vor allem Entitäten, das LLM weist ihnen
vor allem falsche Typen zu.
Ein Kontrollexperiment, in dem die Typbezeichnungen auf CoNLL-2003 und
CleanCoNLL durch bedeutungslose Symbole ersetzt werden, verändert die
Ergebnisse kaum, was die Autoren gegen eine bloße Memorisierung von
Entität-Typ-Paaren anführen.
Zum Ressourcenaufwand enthält die aktuelle Fassung nur einen Vergleich
der Trainingskosten von LLaMA3.1-8B und BERT-MRC auf einer einzelnen GPU,
bei dem Trainingszeit und Spitzenspeicher des LLM moderat höher
ausfallen; Inferenzlatenz und Decoding-Einstellungen werden nicht
berichtet.
Die Arbeit ist zur Veröffentlichung in den Findings der EMNLP 2026
angenommen.

Den Ressourcenaspekt stellen \textcite{obeidat2025llm} in den
Mittelpunkt.
Für fünf biomedizinische Datensätze, darunter ein Social-Media-Korpus zu
Substanzkonsum, vergleichen sie vollständig feinjustierte Encoder mit
rund 340 bis 435~Mio. Parametern (BERT-large, BiomedBERT-large und
DeBERTa-v3-large) mit Mistral-7B-Instruct und Llama-3.1-8B-Instruct, die
mit QLoRA und 4-Bit-Quantisierung angepasst werden; der LoRA-Rang ist so
gewählt, dass die trainierbaren Parameter der LLMs etwa der Größe der
Encoder entsprechen.
Beide Modellklassen geben IOB2-Tags aus und werden strikt auf
Entitätsebene bewertet.
Die LLMs liegen auf vier Datensätzen um 2 bis 8 F$_1$-Punkte vorn, auf
dem fünften knapp (0{,}5 Punkte) hinter dem besten
Encoder~\cite[Tab.~III]{obeidat2025llm}; der Vorsprung
ist bei Entitäten ab drei Tokens am größten und geht vor allem auf eine
höhere Precision zurück.
Dem stehen Inferenzzeiten gegenüber, die je Testbeispiel um das 40- bis
220-Fache über denen des besten Encoders
liegen~\cite[Tab.~IV]{obeidat2025llm}; die LLMs benötigten zudem zwei
statt einer H100-GPU.
Die Autoren führen die Latenz unter anderem darauf zurück, dass das
IOB2-Format für jedes Token ein Tag generieren lässt und lange
Typbezeichnungen in mehrere Tokens zerfallen.

Dass LLMs für tokenbasierte NER ohne Anpassung kaum geeignet sind, zeigen
\textcite{lu2024llmstruggle} am klinischen Korpus RareDis-v1 mit fünf
Typen zu seltenen Krankheiten.
Lokale LLMs mit 7~Mrd. Parametern sowie GPT-3.5 und GPT-4 werden
aufgefordert, IOB2-Tags zu erzeugen.
Im Zero-Shot-Regime erreicht LLaMA-2-7B für den Typ \emph{Rare Disease}
einen F$_1$ von 0{,}149 und GPT-4 einen von 0{,}576.
Nach einem LoRA-Fine-Tuning nähert sich das medizinisch vortrainierte
Llama2-MedTuned-7B nach Darstellung der Autoren dem feinjustierten
BioClinicalBERT an, bleibt aber auf allen Typen darunter; die Autoren
vermuten, dass vollständiges Fine-Tuning den Abstand schließen könnte.
Diese Vermutung relativieren \textcite[Tab.~11]{zhan2026generativeNER}:
In ihren Prüffällen auf CoNLL-2003 und GENIA liegt LoRA um 0{,}14 bis
0{,}48 F$_1$-Punkte über vollständigem Fine-Tuning, was für Qwen3.5 und
MultiNERD jedoch keine Gleichwertigkeit belegt.
Als Gründe für die Schwierigkeiten nennen sie, dass tokenbasierte NER auf
die Repräsentation einzelner Tokens angewiesen ist, dass Decoder im
Pre-Training kein auf Repräsentationen ausgerichtetes Ziel verfolgen und
dass IOB2-Tags für die Modelle ungesehene Symbole sind.
Die Modellgröße variiert die Arbeit nicht.

Einen konstruktiven Umgang mit genau diesem Unterschied zeigen
\textcite{dukic2024lookingright}.
Sie setzen auf Llama-2-7B und Mistral-7B einen klassischen
Klassifikationskopf für Sequence Labeling und entfernen während eines
QLoRA-Fine-Tunings die kausale Aufmerksamkeitsmaske schichtweise, sodass
die Decoder bidirektional auf den Kontext zugreifen können.
Auf CoNLL-2003 steigt der F$_1$ von Llama-2-7B durch die Entfernung der
Maske in allen Schichten von 0{,}744 auf 0{,}908, und beide Decoder
übertreffen RoBERTa-large mit 0{,}900; bei einem Modell mit nur 68~Mio.
Parametern bringt die Maskenentfernung dagegen keinen Gewinn.
Die Decoder werden hier also nicht generativ, sondern wie Encoder
eingesetzt, was die Formatfrage umgeht, aber eine Suche über
Maskierungskonfigurationen erfordert.

Drei der vier Arbeiten konvergieren in einem Punkt: Mit Adaptern
feinjustierte LLMs erreichen auf etablierten Benchmarks das Niveau
feinjustierter Encoder oder übertreffen es
leicht~\cite{zhan2026generativeNER,obeidat2025llm,dukic2024lookingright};
\textcite{lu2024llmstruggle} finden im klinischen Setting mit LoRA
dagegen einen bleibenden Rückstand.
Generativ arbeiten dabei nur \textcite{zhan2026generativeNER} und
\textcite{obeidat2025llm}, während \textcite{dukic2024lookingright} die
Decoder mit Klassifikationskopf einsetzen; den Mehraufwand bei der
Inferenz beziffern allein \textcite{obeidat2025llm}, und dort fällt der
Leistungsgewinn im Verhältnis dazu klein aus.
Zugleich markieren sie die Grenzen des bisherigen Wissens.
Die Datensätze sind allgemeine Benchmarks wie CoNLL-2003 mit vier und
OntoNotes mit 18 Typen oder Fachkorpora; ein genreübergreifender
Datensatz mit 15 Typen wie MultiNERD ist nicht darunter.
Die generativen Ausgabeformate wachsen mit der Eingabelänge: Die
IOB2-Ausgaben bei \textcite{obeidat2025llm} führen jedes Eingabetoken
mit seinem Tag auf, die Inline-Markierungen bei
\textcite{zhan2026generativeNER} wiederholen den Satz; beides erhöht
nach den Befunden aus Abschnitt~\ref{sec:llm_ner_related} die Latenz.
Die Ressourcenerfassung ist unvollständig: \textcite{obeidat2025llm}
messen Inferenzzeiten, aber keinen Speicher, \textcite{zhan2026generativeNER}
berichten nur Trainingskosten, und \textcite{lu2024llmstruggle} sowie
\textcite{dukic2024lookingright} messen keine Ressourcen.
Die untersuchten LLMs liegen durchweg zwischen 3 und 9~Mrd. Parametern,
sodass weder deutlich kleinere noch deutlich größere Decoder im Vergleich
stehen.
Und die Encoder-Referenzwerte stammen teilweise aus anderen Publikationen
statt aus derselben Pipeline.
Diese Punkte bilden den Ausgangspunkt der Forschungslücke in
Abschnitt~\ref{sec:research_gap}.

\section{Einfluss der Modellgröße in verwandten Arbeiten}
\label{sec:model_size_related}

Die in Abschnitt~\ref{sec:model_size} dargestellten Skalierungsgesetze
beschreiben den Pre-Training-Verlust in Abhängigkeit von Parameterzahl,
Datenmenge und Rechenaufwand.
Welche Rolle die Modellgröße für eine nachgelagerte Aufgabe wie NER
spielt, lässt sich daraus nicht ableiten.
Die zuvor vorgestellten Arbeiten erlauben hierzu nur indirekte Aussagen,
weil die Modellgröße fast immer gemeinsam mit anderen Faktoren variiert.

Innerhalb einer LLM-Familie sind Größeneffekte sichtbar, aber klein.
UniversalNER gewinnt beim Übergang vom 7B- zum 13B-Studenten im
Zero-Shot-Regime leicht hinzu, was die Autoren als bessere
Generalisierung größerer Modelle deuten~\cite{zhou2024universalner}.
GoLLIE zeigt von 7 über 13 bis 34~Mrd. Parametern im Mittel steigende
Zero-Shot-Werte, wobei einzelne Datensätze stark, andere gar nicht
profitieren; der Vergleich der berichteten Mittelwerte zeigt zudem, dass
der Gewinn durch die Vergrößerung des Backbones deutlich hinter dem
Gewinn durch das Richtlinientraining zurückbleibt, und auf MultiNERD
liegt das 13B-Modell nicht hinter dem 34B-Modell~\cite{sainz2024gollie}.
\textcite{zhan2026generativeNER} unterscheiden eine leichte und eine
mittlere Größenklasse, ziehen aber keinen Skalierungsschluss.
Zwischen den Modellklassen fallen die Vergleiche noch uneindeutiger aus:
GLiNER und NuNER erreichen mit Encodern weit unterhalb einer Milliarde
Parameter das Niveau von LLMs mit 7 bis 13~Mrd.
Parametern~\cite{zaratiana2024gliner,bogdanov2024nuner}.
Dabei teilen nur GLiNER und UniNER die Pile-NER-Trainingsdaten, NuNER
nutzt GPT-3.5-Annotationen des C4-Korpus, sodass der Größenunterschied
von Architektur, Trainingsziel und Daten überlagert wird.
Bei \textcite{obeidat2025llm} ist die Zahl trainierbarer Parameter über
den LoRA-Rang auf eine ähnliche Größenordnung wie die Encoder abgestimmt
(rund 341 bzw. 353~Mio. Adapterparameter gegenüber 336 bis 435~Mio.),
und dennoch liegen die LLMs vorn.
Dieser Vorsprung dürfte dann eher auf die Kapazität des eingefrorenen
Basismodells als auf die Zahl angepasster Parameter zurückgehen.
Bei Encodern selbst ist die größere Variante nicht durchgehend besser:
In der Vergleichsstudie von \textcite{keraghel2024recent} liegt
DeBERTa-v3-base auf CoNLL-2003 und NCBI Disease knapp vor
DeBERTa-v3-large, auf den meisten anderen Datensätzen ist es umgekehrt.
Und \textcite{dukic2024lookingright} beobachten, dass ihre
Maskenentfernung erst ab einer gewissen Modellgröße wirkt.

Für die vorliegende Arbeit folgt daraus zweierlei.
Erstens wird die Modellgröße innerhalb jeder Familie getrennt betrachtet:
auf der Encoder-Seite mit DeBERTa-v3-base (rund 184~Mio. Parameter) und
DeBERTa-v3-large (rund 434~Mio.), auf der LLM-Seite mit Qwen3.5-0.8B,
-4B und -27B, deren Textmodelle rund 0{,}75, 4{,}21 und 26{,}90~Mrd.
Parameter umfassen (Abschnitt~\ref{sec:groessenklassen}).
Die Spanne reicht damit vom kleinsten Encoder bis zum größten LLM über
mehr als zwei Größenordnungen und deckt auf der LLM-Seite sowohl ein
Modell unterhalb als auch ein Modell oberhalb des in den direkten
Vergleichen der Literatur abgedeckten Größenbereichs ab.
Zweitens bleibt die Auswertung deskriptiv: Mit zwei Encoder- und drei
LLM-Größen, größenabhängig gewählten Hyperparametern und
unterschiedlichen GPU-Typen (Abschnitt~\ref{sec:methodischer-ansatz})
lassen sich Größeneffekte gegen die Streuung über Seeds abgrenzen, aber
keine Skalierungsgesetze schätzen.
Die Frage, ob der Mehraufwand größerer LLMs in einem angemessenen
Verhältnis zum Leistungsgewinn steht, wird deshalb für beide LLM-Regime
je Größenklasse beantwortet.

\section[Forschungslücke und Einordnung der eigenen Arbeit]{Forschungslücke und Einordnung der \mbox{eigenen} Arbeit}
\label{sec:research_gap}

Die vorangegangenen Abschnitte zeigen, dass direkte Vergleiche zwischen
Encodern und generativen LLMs für NER existieren und in ihren
Ergebnissen weitgehend übereinstimmen.
Die Forschungslücke liegt daher nicht im Fehlen solcher Vergleiche,
sondern in ihrer Reichweite.
Tabelle~\ref{tab:related-overview} stellt die behandelten Arbeiten nach
Modellen, Daten, Anpassungsverfahren, Art des Encoder-LLM-Vergleichs und
erfassten Ressourcen zusammen und ordnet die vorliegende Arbeit ein.

\begin{table}[p]
  \centering
  \scriptsize
  \setlength{\tabcolsep}{3pt}
  \renewcommand{\arraystretch}{1.15}
  \caption[Übersicht der behandelten Arbeiten]{Übersicht der behandelten Arbeiten.
  Vergleich Encoder--LLM: „ja“ = beide Modellklassen in derselben Studie
  trainiert bzw. ausgewertet; „berichtete Werte“ = Vergleich mit aus
  anderen Publikationen übernommenen Werten; „--“ = kein Vergleich.
  Ressourcen gemessen: in der Studie erhobene Trainings- oder
  Inferenzgrößen; „--“ = keine Messung. Full-FT: vollständiges
  Fine-Tuning aller Gewichte.}
  \label{tab:related-overview}
  \begin{tabular}{@{}%
      >{\raggedright\arraybackslash}p{2.4cm}%
      >{\raggedright\arraybackslash}p{3.0cm}%
      >{\raggedright\arraybackslash}p{2.9cm}%
      >{\raggedright\arraybackslash}p{2.0cm}%
      >{\raggedright\arraybackslash}p{2.0cm}%
      >{\raggedright\arraybackslash}p{2.2cm}@{}}
    \toprule
    \textbf{Arbeit} & \textbf{Modelle} & \textbf{Daten/Domäne} &
    \textbf{Anpassung} & \textbf{Vergleich Encoder--LLM} &
    \textbf{Ressourcen gemessen} \\
    \midrule
    \textcite{tedeschi2022multinerd} & mBERT + BiLSTM-CRF &
      MultiNERD (10 Spr., 15 Typen); Gold-Benchmarks & Fine-Tuning &
      -- & -- \\
    \textcite{tan2023damo} & XLM-R-large + Retrieval; ChatGPT &
      MultiCoNER II (12 Spr., 33 Typen) & Fine-Tuning; Prompting &
      ja (Validierung) & -- \\
    \textcite{keraghel2024recent} & BERT, XLM-R, DeBERTa-v3, GLiNER;
      GPT-4o & 10 Datensätze, allgemein und Domänen &
      Fine-Tuning; Prompting mit Beispielentitäten & ja & -- \\
    \textcite{zaratiana2024gliner} & DeBERTa-v3 (50--300 Mio.) &
      Pile-NER; CrossNER, MIT & Pre-Training; Zero-Shot &
      berichtete Werte & -- \\
    \textcite{bogdanov2024nuner} & RoBERTa-base; UniNER-7B; GPT-3.5/4 &
      C4-Annotationen; 4 Datensätze & Pre-Training; Few-Shot-FT &
      ja & Kosten der Annotation (unter 500~\$); keine Latenz- oder Speichermessung \\
    \textcite{semin2026spanLabeling} & Qwen3, Mistral, Llama, Gemma
      (8--70 Mrd.) & UNER-Benchmark u.\,a. & Prompting & -- &
      Ausgabetokens \\
    \textcite{wang2025gptner} & GPT-3 & CoNLL, OntoNotes, ACE, GENIA &
      Few-Shot-Prompting (Retrieval mit feinjustiertem Tagger) &
      berichtete Werte & -- \\
    \textcite{xie2023zeroshot} & GPT-3.5 & 7 Benchmarks (EN, ZH) &
      Zero-Shot-Prompting & berichtete Werte & -- \\
    \textcite{ashok2023promptner} & GPT-4, GPT-3.5, T5 &
      CoNLL, GENIA, FewNERD, CrossNER & Few-Shot-Prompting &
      berichtete Werte & -- \\
    \textcite{li2023codeie} & Codex, GPT-3 & CoNLL, ACE;
      RE-Datensätze & Few-Shot-Prompting & berichtete Werte (UIE) & -- \\
    \textcite{zhou2024universalner} & LLaMA 7B/13B &
      43 Datensätze, 9 Domänen (inkl. MultiNERD-Variante) &
      Destillation; Fine-Tuning & berichtete Werte & -- \\
    \textcite{sainz2024gollie} & Code-LLaMA 7/13/34B &
      IE-Datensätze; Zero-Shot u.\,a. CrossNER, MultiNERD & QLoRA & -- &
      Trainingszeit, CO$_2$ je Größe (keine Inferenz) \\
    \textcite{ding2024gner} & Flan-T5 (bis 11 Mrd.), LLaMA-7B &
      Pile-NER; Zero-Shot CrossNER, MIT & Fine-Tuning (Label je Wort) &
      berichtete Werte (GLiNER) & Durchsatz; GPU-Stunden des Trainings \\
    \textcite{zhan2026generativeNER} & acht LLMs (3--9 Mrd.);
      BERT-MRC u.\,a. & CoNLL, OntoNotes, ACE2005, GENIA & LoRA &
      teilweise (BERT-MRC selbst trainiert) & Trainingskosten (keine Inferenz) \\
    \textcite{obeidat2025llm} & BERT-large, BiomedBERT,
      DeBERTa-v3-large; Mistral-7B, Llama-3.1-8B &
      5 biomedizinische Datensätze & Full-FT; QLoRA & ja &
      Inferenzzeit, GPU-Anzahl \\
    \textcite{lu2024llmstruggle} & 7B-LLMs; GPT-3.5/4; BioClinicalBERT &
      RareDis-v1 (klinisch) & Prompting; LoRA & ja & -- \\
    \textcite{dukic2024lookingright} & Llama-2-7B, Mistral-7B; RoBERTa &
      CoNLL, Rest14, ACE05 & QLoRA mit Maskenentfernung & ja & -- \\
    \textcite{lu2024padellm} & Llama-2-7B, Baichuan2-7B &
      CoNLL, ACE, GENIA; chin. Datensätze & Fine-Tuning & -- &
      Inferenzlatenz \\
    \midrule
    Diese Arbeit & DeBERTa-v3-base/-large; Qwen3.5-0.8B/\allowbreak-4B/\allowbreak-27B &
      MultiNERD (EN, 15 Typen); CrossNER-Transfer &
      Full-FT; Zero-Shot; QLoRA & ja (gemeinsame Pipeline) &
      Trainingsdauer, Latenz, GPU-Speicher \\
    \bottomrule
  \end{tabular}
\end{table}

Die erste Lücke betrifft Datenbasis und Pipeline, die ersten beiden der
vier Aspekte aus Abschnitt~\ref{sec:problem_ziel}.
Die direkten Vergleiche verwenden allgemeine Benchmarks wie CoNLL-2003
mit vier und OntoNotes mit 18 Typen aus mehreren Genres wie
Nachrichten, Rundfunk, Gesprächen und Webtexten~\cite{weischedel2013ontonotes}
oder Fachkorpora wie die biomedizinischen Datensätze bei
\textcite{obeidat2025llm}.
MultiCoNER~II bei \textcite{tan2023damo} ist mit 33 Typen zwar
feingranular, besteht aber aus kurzen, kontextarmen Sätzen und stellt
den Encodern nur ein per Prompting genutztes ChatGPT gegenüber.
Publizierte MultiNERD-Werte beruhen auf übernommenen Baselines oder
eigener Aufbereitung: \textcite{zhou2024universalner} übernehmen die
BERT-Werte aus InstructUIE auf dessen Datensatzvariante,
\textcite{zaratiana2024gliner} werten GLiNER auf der
InstructUIE-Aufbereitung mit 10\,000 Testbeispielen aus, und bei
\textcite{sainz2024gollie} bleibt das Verhältnis der 32\,908
Testbeispiele zu den offiziellen Splits offen.
Eine gemeinsame Auswertung beider Modellfamilien auf den offiziellen
englischen MultiNERD-Splits mit einer in derselben Pipeline trainierten
Encoder-Referenz fehlt damit; zudem beruhen die Encoder-Referenzwerte
der direkten Vergleiche teilweise auf übernommenen Baselines.
Die vorliegende Arbeit setzt hier an, indem sie DeBERTa-v3-base und
-large sowie Qwen3.5 in drei Größen auf den offiziellen englischen
MultiNERD-Splits in einer gemeinsamen Pipeline trainiert und auswertet,
mit derselben \texttt{seqeval}-Implementierung und derselben
Silberreferenz, und dabei Trainingsdauer, Inferenzlatenz und
GPU-Spitzenspeicher der Inferenz nach einem einheitlichen Protokoll
erfasst (Abschnitt~\ref{sec:effizienzmetriken}).
Mit der Indexing-Ausgabe wird zudem ein Format verwendet, das den
Eingabetext nicht reproduziert und dessen Formtreue eigens gemessen wird.
Diese Lücke betrifft vor allem FF1, den Vergleich von
Erkennungsleistung und Formtreue in einer gemeinsamen Pipeline, sowie
FF3, die Ableitung von Empfehlungen für die Auswahl geeigneter
NER-Systeme unter Speicher- und Latenzbeschränkungen
(Abschnitt~\ref{sec:problem_ziel}).

Die zweite Lücke betrifft Ressourcenerfassung und Modellgröße, den
dritten und vierten Aspekt.
Leistung und Ressourcen werden in den direkten Vergleichen nicht
gemeinsam erfasst: Inferenzzeiten ohne Speicher, Trainingskosten ohne
Inferenz oder gar keine Messung.
Zudem betrachten die direkten Vergleiche LLMs einer einzigen
Größenklasse; Größenstaffelungen finden sich nur innerhalb von
LLM-Familien, bei GoLLIE mit Trainingsressourcen je Größe, aber ohne
gemeinsamen Encodervergleich mit
Inferenzmessung~\cite{zhou2024universalner,sainz2024gollie}, während
Encoder-Studien die Größe nur innerhalb ihrer Klasse
variieren~\cite{keraghel2024recent}.
Keine der Arbeiten in Tabelle~\ref{tab:related-overview} untersucht eine
Größenstaffelung auf beiden Seiten mit gemeinsamer Leistungs- und
Ressourcenerfassung.
Die vorliegende Arbeit vergleicht drei Größen einer LLM-Familie zwischen
0{,}8 und 27~Mrd. Parametern in zwei Regimen, Zero-Shot und QLoRA, mit
zwei Encoder-Größen derselben Architektur und erfasst für jede
Konfiguration Leistung, Formtreue und Ressourcen.
Sie adressiert damit FF2, die Veränderung von Erkennungsleistung und
Ressourcenbedarf mit der Modellgröße, und liefert die Grundlage für die
größenbezogenen Empfehlungen in FF3.

Zwei Einschränkungen sind offen zu benennen.
Die erste betrifft eine mögliche Datenkontamination.
MultiNERD besteht aus Wikipedia- und Wikinews-Artikeln, die mit hoher
Wahrscheinlichkeit Teil der Pre-Training-Korpora der verwendeten LLMs
sind; die Zusammensetzung der Qwen3.5-Trainingsdaten ist nicht
vollständig offengelegt, sodass sich dies weder belegen noch ausschließen
lässt.
\textcite{zhan2026generativeNER} weisen auf bekannte Lecks von
Benchmark-Daten in LLM-Trainingsdaten hin und ziehen deshalb mit
CleanCoNLL einen Datensatz heran, dessen bereinigte Form kaum im
Pre-Training enthalten sein dürfte; ihr Befund, dass die Ergebnisse nicht
von den Typbezeichnungen abhängen, spricht gegen eine Memorisierung von
Entität-Typ-Paaren, gilt aber für CoNLL-2003 und CleanCoNLL.
Für MultiNERD bleibt die Frage offen.
Zu bedenken ist dabei, dass die Referenz ein automatisch erzeugter Silver
Standard ist: Ein Modell, das den Wikipedia-Text kennt, kennt damit noch
nicht die hier verwendeten Labels, auch wenn der Datensatz selbst seit
2022 öffentlich verfügbar ist.
Im QLoRA-Regime lernen alle Modelle das Typinventar ohnehin aus den
Trainingsdaten; im Zero-Shot-Regime könnte Vorwissen über die Artikel die
Ergebnisse dagegen begünstigen.
Die Arbeit kann diese Frage nicht klären; sie wird bei der Interpretation
der Zero-Shot-Ergebnisse berücksichtigt.

Die zweite Einschränkung betrifft die Vergleichsbedingungen.
Anders als eine idealisierte Studie hält die vorliegende Arbeit nicht
alle Bedingungen konstant: Die Konfigurationen liefen auf drei
GPU-Typen, die Hyperparameter wurden je Modellgröße gewählt, und Prompt,
Decoding und numerische Präzision unterscheiden sich zwischen Zero-Shot-
und QLoRA-Läufen (Abschnitt~\ref{sec:validitaet}).
Konstant sind Datenbasis, Splits, Evaluationsprotokoll und
Bibliotheksversionen.
Die Arbeit vergleicht deshalb konkrete Konfigurationen und bestimmt keine
isolierten kausalen Effekte einzelner Faktoren; Effizienzvergleiche
werden vorrangig innerhalb gleicher Hardware- und Präzisionsbedingungen
interpretiert.
Ähnliche Einschränkungen finden sich auch in den direkten Vergleichen
der Literatur: \textcite{obeidat2025llm} betreiben die LLMs auf zwei und
die Encoder auf einer GPU, \textcite{zhan2026generativeNER} berichten
Trainingskosten, aber keine Inferenzbedingungen; hier werden die
Abweichungen explizit ausgewiesen.
Der CrossNER-Ergänzungsversuch schließlich erweitert den Vergleich um
einen reinen Transfertest der jeweils größten Modelle beider Familien auf
die Grobtypen PER, ORG und LOC; er ist, wie in
Abschnitt~\ref{sec:llm_ner_related} dargelegt, nicht mit den Few-Shot-
und Richtlinienansätzen der Literatur vergleichbar und dient der
Einordnung, nicht dem Benchmark-Vergleich.

Aus der Literatur folgen drei Erwartungen als Lesehilfe, die
Kapitel~\ref{chap:fazit} je Forschungsfrage mit den Befunden abgleicht;
ein Hypothesentest ist damit nicht beabsichtigt.
E1 (zu FF1): Ein mit LoRA oder QLoRA angepasstes LLM erreicht das
Niveau eines vollständig feinjustierten Encoders oder übertrifft es
leicht~\cite{zhan2026generativeNER,obeidat2025llm}.
E2 (zu FF2 und FF3): Die Inferenz generativer LLMs ist um ein Vielfaches
langsamer als die eines Encoders; die Literatur nennt Faktoren von etwa
40 bis 220~\cite[Tab.~IV]{obeidat2025llm}.
E3 (zu FF1): Ohne Anpassung der Gewichte, also im Zero-Shot-Regime,
bleiben per Prompting genutzte LLMs deutlich hinter feinjustierten
Modellen zurück~\cite{keraghel2024recent,tan2023damo}.
